\documentclass[UTF8,a4paper,zihao=-4,fontset=none]{ctexart}
\usepackage{orchid-report}

\newcommand{\reportauthor}{}
\newcommand{\reportid}{}
\newcommand{\reportclass}{}
\newcommand{\reportprogram}{}
\newcommand{\reportdate}{2026 年 11 月 21 日}
\InputIfFileExists{cover.local.tex}{}{}
\hypersetup{pdftitle={Orchid 从零训练的五言诗生成与韵律控制},pdfauthor={\reportauthor}}

\begin{document}
\begin{titlepage}
\thispagestyle{empty}
\centering
\setlength{\parindent}{0pt}
\vspace*{16mm}
{\color{pine}\large 课程项目报告\par}
\vspace{13mm}
{\fontsize{67}{74}\selectfont\color{ink} Orchid\par}
\vspace{4mm}
{\fontsize{23}{32}\selectfont\bfseries\color{pine} 从零训练的五言诗生成与韵律控制\par}
\vspace{6mm}
{\color{line}\rule{.56\textwidth}{.7pt}\par}
\vspace{10mm}
\noindent\begin{minipage}{.59\textwidth}
  \begin{poembox}
    \centering
    {\versefont\Large
    霜葉滿空落，\\[2pt]
    悲心無處{\color{clay}尋}。\\[2pt]
    孤舟寒夜水，\\[2pt]
    獨自照秋{\color{clay}林}。\par}
    \vspace{2mm}
    {\small\color{muted}「霜、孤舟」实测输出\par
    平水韵下平声十二侵\par}
  \end{poembox}
\end{minipage}\par
\vfill
\ifdefempty{\reportauthor}{}{%
  {\large\bfseries\reportauthor\enspace \reportid\par}
  \vspace{1.6mm}
  {\color{muted}\reportclass\par
  \reportprogram\par}
  \vspace{5mm}
}
{\color{pine}\reportdate\par}
\vspace*{10mm}
\end{titlepage}

\pagenumbering{Roman}
\section*{摘要}
\addcontentsline{toc}{section}{摘要}

Orchid 是一个用公开唐宋诗从随机初始化开始训练的字符级 GPT。项目任务是在单张消费级显卡上完成语料清洗、分词、模型训练、专用微调和推理部署。输入为一至三个短关键词，目标为四句、每句五字的诗。工程采用诗作级分组划分，使完全重复、可检测近重复和长诗内嵌片段尽量留在同一集合。唐宋扩容沿用唐诗字表与冻结留出集，得到 296,821 首训练诗作。模型从约 14M 扩至 92M，训练代码保留因果多头注意力、残差连接、绝对位置嵌入和逐标记语言模型目标。

在同一唐诗语料及约 2,849 万有效训练目标下，14M、42M、92M 的验证负对数似然分别为 4.9540、4.8151、4.6831。唐宋 92M 的最佳验证值达到 3.8805。随后用仅对诗文和结束符计损失的监督配对训练条件创作，再以来自同一初始化的普通续训组检验同韵样本的作用。最终权重在 50 道冻结新题、每题三次采样的自由模式下有 148/150 首符合四句五字，92/150 首可确认偶句异字同韵，17/150 首同时满足句式、押韵和原词覆盖。默认约束解码三项均达到 150/150。严格逐字平仄常格及相应押韵的联合判定有 1 首通过，说明偶句同韵覆盖了格律要求的一部分。

补充实验固定权重、候选池和采样种子，对比两种候选排序。新加入的隔位重字惩罚在 36 道独立编排的新题上将有该类重复的输出从 13/108 降到 9/108。两份基于独立 AI 上下文的 v2 盲评均给出较高的切题增益，可读性变化方向不一致。项目提供同题对照 HTML、离线样例、常驻本地模型接口、可复跑的统计脚本与数据指纹。

\vspace{.45em}
\noindent\textbf{关键词}\quad 字符级语言模型、五言诗、条件生成、约束解码、平水韵

\clearpage
\begingroup
\small\linespread{1.07}\selectfont
\setlength{\parskip}{0pt}
\tableofcontents
\endgroup
\clearpage
\pagenumbering{arabic}

\section{项目任务与技术路线}
\label{sec:task}

输入「霜、孤舟」时，系统要在短小的篇幅内安置两个意象，写成四句五字。课程要求覆盖从模型搭建、预训练、微调到部署的完整过程。选择古诗是因为字数、句数、韵脚和字面用词可以复核，切题、可读性与意境又需要读原诗。这样，同一次实验可以检验程序是否达成任务约定，也能观察自动指标以外的生成效果。

Orchid 实现三个规模的字符级 GPT。14M 用于跑通基线和检查训练程序，42M 与 92M 用于同预算对照。最终系统以 92M 唐宋模型为起点，经监督微调和同韵诗续训，再在推理时按需限制字数、偶句韵部与用词。系统提供自由模式，使训练所得能力与推理时额外保障的规则可以在同一权重上分开测量。图~\ref{fig:workflow} 给出完整流程和验证所在的位置。

\reportfigure{workflow.pdf}{项目流程。公开语料先分组再派生分词窗口及监督样本。冻结唐诗验证集用于选模，独立新题用于报告生成表现。}{fig:workflow}

\begin{table}[!htbp]
  \centering
  \caption{报告覆盖 \texttt{Problem.md} 的项目评分要求}
  \label{tab:rubric}
  \begin{tabularx}{\linewidth}{@{}p{.23\linewidth}p{.10\linewidth}Y@{}}
    \toprule
    评分项 & 占比 & 对应内容与证据\\
    \midrule
    知识与原理 & 35\% & 自回归、因果注意力、位置嵌入、参数量推导、分词与监督目标，见第~\ref{sec:model}、\ref{sec:corpus} 和 \ref{sec:sft} 节。\\
    模型实现 & 20\% & 可训练的解码器、三种配置、同预算实验、单卡容量检查与部署，见第~\ref{sec:model}、\ref{sec:pretrain} 和 \ref{sec:deployment} 节。\\
    预训练 & 20\% & 固定快照、训练窗、随机初始化、优化日志与续训、学习曲线和原诗续写，见第~\ref{sec:corpus} 和 \ref{sec:pretrain} 节。\\
    微调 & 20\% & 监督标签、掩码、选模、韵律对照与生成测试，见第~\ref{sec:sft} 至 \ref{sec:evaluation} 节。\\
    成果与体会 & 5\% & 交付范围、失败样例、限制、复现记录与个人项目体会，见第~\ref{sec:deployment} 和 \ref{sec:conclusion} 节。\\
    \bottomrule
  \end{tabularx}
\end{table}

\subsection{已有工作与本项目取舍}

Transformer 以注意力为核心，GPT-2 展示了单向语言模型的自回归学习路线~\cite{vaswani2017attention,radford2019language}。中文古诗自动生成早期已有循环网络方案~\cite{zhang-lapata-2014-chinese}。九歌将押韵和诗体检查接入生成系统~\cite{guo-etal-2019-jiuge}，CharPoet 研究逐字生成对格式控制的作用~\cite{yu-etal-2024-charpoet}。本项目使用随机初始化的字符模型和固定唐宋诗料。论文结果用来解释设计选择，正文中的各组数值均来自 Orchid 自己的协议和实测。

\section{模型结构与学习目标}
\label{sec:model}

\subsection{自回归概率与目标函数}

设诗文经分词后的序列为 $x_0,\ldots,x_T$。训练时将输入与标签错开一个位置。第 $t$ 个位置可看见 $x_{\le t}$，输出对 $x_{t+1}$ 的未归一化分数。令 $\mathcal I$ 为参与计损失的位置集合，预训练的字词预测与监督微调可统一写成
\begin{equation}
    \mathcal L =
    -\frac{1}{|\mathcal I|}\sum_{t\in\mathcal I}
    \log p_\theta(x_{t+1}\mid x_{\le t}),
    \qquad \mathrm{PPL}=\exp(\mathcal L).
    \label{eq:nll}
\end{equation}
代码用交叉熵直接接收 logits，将批次的损失和按有效目标数汇总。有效目标数对不同长度的样本和最后一个小批次都能正确计权。训练借助已知前缀并行预测各位置，运行时则每次从最后位置选一个标记并接回序列，遇结束符或达到预算后停止。监督阶段和预训练阶段的目标集合不同，比较 NLL 时始终在同一口径和相同留出集内进行。

\subsection{嵌入、因果注意力与前馈块}

每个位置将字符嵌入与可学习的绝对位置嵌入相加。对输入 $X\in\mathbb R^{T\times d}$，每个注意力头分别由线性投影产生 $Q,K,V$。预归一化块以 LayerNorm 后的表示计算
\begin{equation}
  \operatorname{Attn}(Q,K,V)
   = \operatorname{softmax}\left(\frac{QK^{\mathsf T}}{\sqrt{d_h}}+M\right)V,
  \qquad
  M_{ij}=
    \begin{cases}
      0,&j\le i,\\
      -\infty,&j>i.
    \end{cases}
\end{equation}
其中 $d_h=d/H$，用 $\sqrt{d_h}$ 缩放点积，可以让头宽变化时的注意力分数保持在较适合 softmax 的量级。掩码让位置 $i$ 只从自己和此前位置取值。将各头结果拼接并投影后，与块输入做一次残差加法。第二个子层使用 $d\to4d\to d$ 的 GELU 前馈网络，并进行另一次残差加法。末端归一化后，用词嵌入矩阵的转置生成词表 logits。

\reportfigure{transformer.pdf}{92M 模型结构和因果掩码。左侧的注意力与前馈组成一个块，重复 12 次。右侧绿色单元允许当前位置访问对应键。}{fig:transformer}

实现提供 PyTorch 的 scaled dot-product attention 与显式掩码路径，并在评估模式下关闭 dropout。测试检查两条注意力路径的输出差异、未来 token 对当前 logits 的隔离，以及权重保存、加载和训练续跑。共用输入嵌入与输出头的权重后，含偏置的参数量为
\begin{equation}
    N=d(V+C)+L(12d^2+13d)+2d,
\end{equation}
其中 $V$ 是词表大小，$C$ 是最大上下文，$L$ 是块数。块内的注意力和前馈分别贡献 $4d^2$ 与 $8d^2$ 个矩阵参数，两组 LayerNorm 及各线性层偏置贡献 $13d$ 个参数。模型共享 $V=8{,}582$、$C=256$、头宽 64 与 0.1 的训练 dropout。代码在实例化后逐张量计数核对配置。

\begin{table}[!htbp]
\centering
\caption{Orchid 的三种模型配置。实际参数数包含嵌入与归一化参数。}
\label{tab:model}
\begin{tabular}{@{}lrrrrr@{}}
\toprule
版本 & 块数 & 隐藏宽度 & 注意力头 & 前馈宽度 & 实际参数\\
\midrule
14M & 6 & 384 & 6 & 1,536 & 14,041,344\\
42M & 12 & 512 & 8 & 2,048 & 42,354,688\\
92M & 12 & 768 & 12 & 3,072 & 91,843,584\\
\bottomrule
\end{tabular}
\end{table}

\section{语料、冻结划分与字符词表}
\label{sec:corpus}

\subsection{唐诗基线与宋诗扩展}

唐宋正文来自 \texttt{chinese-poetry} 仓库的同一固定提交，唐诗 58 个分片，宋诗 255 个分片~\cite{chinesepoetry}。预处理从 \code{paragraphs} 提取正文，做 Unicode NFC 归一化，凭句读恢复诗句，将标点规范为逗号与句号交替。原始文字和来源索引保留供复查。无句子、缺字标记和混入非汉字内容的记录被登记。标题及作者用于追溯来源。监督目标必须本身可解析为四句五字，长诗只用于预训练与重合检查。

唐诗 57,607 条原始记录有 51,945 条通过初筛，按正文去重后为 51,138 首。按诗作分组，在相同或可检测的近重复关系上建立连边，再做约 90/5/5 的固定划分。训练、验证、测试分别有 46,031、2,557、2,550 首。近重复规则检查长诗内嵌的四句、四句诗三句对齐及有限编辑距离，让相关文本尽量同组。训练和评估窗口继承原有划分。这个边界特别重要，逐行随机分配会把同一首诗的前后句分到不同集合。

\reportfigure{corpus.pdf}{语料清洗与训练集扩容。右侧新增的宋诗只进训练集，唐诗验证与测试的诗文、标记与监督样本均按指纹核验。}{fig:corpus}

扩容时对宋诗重复执行正文清洗，筛除唐诗同文和与留出诗符合近重复规则的候选。宋诗中清洗、去重及泄漏筛查后保留 250,790 首，新训练集合达 296,821 首，词流含 18,451,787 个标记。宋诗加入后仍使用唐诗训练集学得的字表，以维持模型结构与留出集口径；扩展训练正文中的未登录 token 约占 0.374\%。该规则可以追踪已知重合，异文和典故借用仍可能穿过规则边界。

\subsection{字符级分词与两类训练样本}

字符词表仅由唐诗训练集拟合，固定顺序为保留标记加按字符排序的正文与标点。保留标记包括序列开头、结尾、填充、未知字，以及专用诗体和关键词分隔符。以字符为单位便于按汉字数施加规则，第 2、4 句的位置也可以落实到确定的生成步。代价是古诗很短时仍需逐字解码，生僻字和简繁异体须受词表覆盖约束。界面保留训练采用的字形，并对未登录字符返回明确提示。

预训练将分好的诗文流按上下文窗切片，每个位置预测其后标记。完整窗有 256 个位置，末尾填充位置的标签屏蔽。监督阶段从确认为五言四句的诗作中构造「关键词 $\to$ 诗文」对。唐诗与唐宋版训练监督集分别为 2,929、12,394 对，两版共有同一组 166 条验证和 157 条测试配对。与写作者给出的题意相比，训练标签只是可定位的弱监督。独立诗题、语义评分与显式约束分别检查条件创作的三个侧面。

\section{单卡预训练与规模实验}
\label{sec:pretrain}

\subsection{训练程序和资源配置}

预训练均从随机权重开始，用 RTX 4080 SUPER 在 Ubuntu 22.04、Python 3.11.16 和 PyTorch 2.14.1+cu126 环境运行。正式训练开启 BF16 自动混合精度，保留 FP32 主参数。优化使用 AdamW~\cite{loshchilov2019decoupled}，$\beta_1=0.9$、$\beta_2=0.95$、$\epsilon=10^{-8}$，二维及以上权重应用权重衰减。梯度全局范数限制为 1.0，学习率先预热再按余弦变化。唐诗等预算组用微批 32、累积 4，完整优化步的目标数为 $32\times4\times256=32{,}768$。计算每个累积窗口的有效目标总数以后再归一化梯度，可避免短窗和末尾批次改变优化权重。

训练程序在检查点保存模型参数、优化器、调度器、随机数状态、数据游标及数据与代码指纹。续训从游标定位到同一轮尚未处理的窗口。实际 14M 基线在第 87 步落盘，然后恢复到第 870 步。唐宋 92M 的长程训练也从第 1,974 步继续，验证日志和最优权重可以对应到累计有效目标数。这使一次长训练可以分多次运行，统计仍使用统一口径。

\subsection{固定训练预算下扩大模型}

14M、42M 和 92M 使用同一唐诗语料、同一字表、10 轮及约 2,849 万有效目标。14M 的验证 NLL 从 9.1361 降至 4.9540，10 轮优化计算约 58.9 秒，验证、采样和写盘单独计时。在相同训练预算下，42M、92M 的验证 NLL 进一步降至 4.8151 和 4.6831。扩容效果见图~\ref{fig:training}A。单卡短窗口测试的吞吐与显存见表~\ref{tab:capacity}，这些数据用于选择批量。

\begin{table}[!htbp]
\centering
\caption{固定唐诗语料与微批 32 的容量测试。显存为 PyTorch 峰值保留量。}
\label{tab:capacity}
\begin{tabular}{@{}lrrr@{}}
\toprule
模型 & 有效目标吞吐（token/s） & 峰值保留显存（GiB） & 10 轮验证 NLL\\
\midrule
14M & 456,778 & 1.88 & 4.9540\\
42M & 185,261 & 3.20 & 4.8151\\
92M & 101,098 & 4.97 & 4.6831\\
\bottomrule
\end{tabular}
\end{table}

预训练阶段已经学会生成古诗风格的短句。14M 的一次固定前缀续写为「{\versefont 春風自有情，獨坐故人歸。}」，随后生成结束符。内容和题意有一定联系，而输出只有两句。这项直接观察说明了预训练能提供字句基础，也给监督阶段提出了明确目标，即让输入词与完整四句建立稳定对应。

\subsection{延长训练与扩大唯一文本量}

只用唐诗的 92M 长程配置上限为 26 轮，实跑 16 轮、约 4,558 万有效目标。最佳验证 NLL 4.2995 出现在第 13 轮。此后三轮未改善并触发早停。复用相同的小语料能继续降低损失，但其收益会趋缓。与之配对分析的唐宋版扩大了唯一训练文本，并保持相同 92M 结构、字表与唐诗留出集。

\reportfigure{training.pdf}{预训练验证 NLL 随累计有效目标的变化。A 固定唐诗语料及目标数。B 比较唐诗长程与唐宋扩展，空心点标记各自最佳验证值。两条长程曲线的语料规模和学习率调度同时变化。}{fig:training}

唐宋版按 16 轮上限、每半轮验证训练。第 9 轮、第 5,076 步达到 NLL 3.8805，PPL 为 48.45，此时累计 166,066,074 个目标。训练于第 10.5 轮早停，处理 193,758,254 个目标，优化计算约 30.7 分钟。结合两条曲线可以看到，增加诗料和训练预算有益于同一留出集的语言建模。由于学习率计划和数据规模一起变化，这一结果只能归因于整套扩容配置。92M 微批从 32 提到 64 时，峰值保留显存升至 8.72 GiB，而短测吞吐从约 10.1 万降到约 9.8 万目标每秒。选择微批需要实测吞吐与稳定性。

\section{条件微调与韵律续训}
\label{sec:sft}

\subsection{监督标签与移位掩码}

公开语料缺少每首诗的人工关键词标签。先在训练集上统计意象词白名单，保留见于至少两首训练诗的词，排除停词、长词内包含的短词及跨句匹配。在每首合格诗中按固定种子选择一至三个词并打乱顺序，训练目标仍是原作的完整诗文。唐宋微调训练共有 12,394 对，纳入监督前剔除 776 条目标正文含未登录字符的匹配诗。这种弱监督能在一般的古诗中构造大量配对，但人类读题时常用的反衬、暗示和同义词不会直接进入字面标签。

输入前缀依次提供 \code{<bos>}、\code{<wujue>}、\code{<kw>}、关键词及 \code{<sep>}。训练时，前缀及填充的目标位置记为 \code{-100}。由于标签向前移一位，分隔符所在位置预测第一字，所以从该位置开始计损失，最后一个标点预测结束符。若把该位置也屏蔽，模型就失去了「看到条件后选首字」的直接反馈。图~\ref{fig:sftmask} 对应一条真实监督样本。

\reportfigure{supervision.pdf}{监督配对的输入和下一标记标签。长条件与长正文略去一部分，索引与目标来自冻结的掩码样例。}{fig:sftmask}

\subsection{从基线调通微调到唐宋版选模}

14M 从第 870 步的最佳预训练权重开始监督训练。开发阶段先比较峰值学习率 $3\times10^{-5}$ 与 $3\times10^{-4}$。在相同验证子集上，后者的末轮句式和全词联合达标由 3/30 升到 13/30，而损失仅有轻微反向变化。正式三种子训练共用预训练初始化，每个种子 5 轮、460 步及 366,125 个有效目标。种子 2027 为该阶段的演示权重，以便在保留的 14M 对照和盲评中追溯。

唐宋 92M 从第 9 轮最佳预训练权重开始，仍用三个种子分别创建优化器。扩大配对后采用 $1.5\times10^{-4}$ 的峰值学习率和最多 25 轮的调度，每个种子因验证早停实际运行 4 轮、1,552 步、1,239,400 个目标。三个最低监督验证 NLL 均出现在第 1 轮，分别为 3.6487、3.6600 和 3.6871，微调前的同口径值为 5.5932。先以各自最低 NLL 定种子内检查点，再按验证诗题的字数、用词及训练重合条件选择交付种子。90 道验证题每题采样三次，三个种子无整首或长诗片段命中的联合成功数分别为 66、79、95，因此选出 2028。选择方案在接触冻结测试结果前完成。

\subsection{同韵样本的匹配续训}

为改善自由生成中的押韵，从 92M 同一个最佳监督权重开始，对照抽取 4,096 条普通训练配对，处理组抽取 4,096 条偶句异字、同一平水韵平声部的配对~\cite{pingshui}。两组在词数上匹配，一词、两词、三词样本量分别为 2,028、1,342 和 726。两组用同一个训练种子、相同优化参数，各训练 6 轮、384 步、614,400 个目标。普通组和同韵组的数据取自原训练集合。在 90 道验证题各采 3 次的自由生成中，原权重、普通组、同韵组的异字同韵数分别为 117、119 和 188。

\reportfigure{finetune.pdf}{监督微调和韵律续训。A 为唐宋 SFT 三个种子的第 1 至第 4 轮验证 NLL。B 固定采样条件，三组的分母均为 270。}{fig:sft}

同韵组的 188/270 比原权重的 117/270 多 71 首。三项同时达标由 42/270 到 80/270，训练集整首匹配仍为 0，因此按事先约定切换到同韵组末权重。这是一组同训练种子的配对实验，对稳定性的估计有限。它支持筛选可靠韵脚样本可以改善当前模型的自由押韵，推广到其他语料或韵部仍需更多重复。

\section{生成规则、候选排序与评价}
\label{sec:evaluation}

\subsection{从自由采样到偶句异字同韵}

自由模式按温度 0.8、top-k 为 50、top-p 为 1 的设置抽样，最多生成 64 个新标记。有限输出空间可以在取样前通过词表掩码约束。设前缀为 $h$，第 $t$ 步允许集合为 $\mathcal A_t(h)$，则将允许集合外的 logits 设为 $-\infty$，再对允许集合抽样。固定句式在每五个汉字后生成规定标点，末句标点后只允许结束符。于是读者看到的原文保持四句和二十个汉字。

押韵模式用固定的平水韵表建立韵脚搭档。第二句末在只有一个可确认平声韵部且有搭档的字中抽样，第四句末只允许同部的其他字。多读、多部的韵脚在自动诊断中保留为待判。要求全部关键词时，采样为各词选择不同诗句及句内落点，并让偶句保留韵脚位。模型可在预设落点之前自然生成目标词。默认比较四个候选，以原模型给完整输出的平均对数似然减去整句和双字重复惩罚进行排序。以上机制只规定偶句同韵和所需用词，首句入韵和逐位平仄通过后续诊断评估。

\subsection{冻结主测试与规则归因}

主测试在 50 道预先登记的新题上运行，词数为一、二、三的题目分别有 16、17、17 道。每题使用 11、22、33 三个采样种子，四种模式使用同一权重。统计四句五字、偶句异字同韵、原词全覆盖及三项同时达标。原词匹配要求输入词完整出现在同一句，便于检查生成约束。语义切题另用量表。独立目录中的复测对照了全部 600 条输出，原文和自动判定均与封存版一致。

\reportfigure{rules.pdf}{四种解码方式的规则对照。每种模式计 150 条。只限字数及押韵时全词覆盖与自由模式相同，加入关键词约束后才出现明显变化。}{fig:rules}

自由生成已有 148/150 首符合四句五字，92/150 首确认偶句异字同韵，三项同时达标 17/150。只限长度后，四句五字升到 150/150，三项联合仍是 17/150。加押韵后，偶句异字同韵升到 150/150，三项联合到 25/150。再加入用词约束及四候选排序，四项均为 150/150。自由模式按词数分层，一词为 25/48 全词命中，两词、三词各 0/51。默认模式分别达到 48/48、51/51、51/51。这一差别指向模型主动用全多个词的能力，而非句长与韵脚规则已经达到的上限。

在默认模式的 150 首中，完整生成文本互异，覆盖 24 个可确认平声韵部。按冻结训练集查询整首、长诗连续片段及近重复，三项均为 0/150。对严格逐字平仄常格及同格押韵另行核验时，1 首合、146 首不合、3 首待判。此口径使用四套固定模板，历史变格、多音字和首句入韵存在需据诗句判断的情形。八首已归档的独立 AI 格律复核供个案核对。发布的文字标签使用「五言四句 · 偶句押韵」，并保留平仄诊断结果。

\subsection{两份独立 AI 盲评}

最新 92M 的主观补测从冻结题抽取 30 道，每种词数 10 道，种子固定为 11。每题给出自由与默认输出，两位 AI 评阅者各读 60 首且只接收本人的匿名表与 v2 量表。题序和组内次序独立随机化，同题比较保留配对关系。早期 14M 归档包含一份真人 v1 评分和一份同版 AI 参照，其模型、量表和输出均有独立记录。表~\ref{tab:rubricv2} 是本轮使用的指标。

\begin{table}[!htbp]
\centering
\caption{主观评分定义。分值 1 至 5，NA 表示该项不适用。}
\label{tab:rubricv2}
\begin{tabularx}{\linewidth}{@{}p{.17\linewidth}Y p{.24\linewidth}@{}}
\toprule
维度 & 判断依据 & 适用范围\\
\midrule
切题 & 是否在语义上回应题意，接纳同义、转述和可由诗句推得的暗示。 & 全部输出\\
可读性 & 读者能否理解用词和句意，接纳合理的古汉语省略、倒装。 & 全部输出\\
连贯性 & 四句场景、时间、叙述视角或情绪是否能衔接。 & 可解析的四句\\
意境 & 画面、情意和余味是否由诗句自然形成。 & 可读性和连贯性均至少 3 分\\
\bottomrule
\end{tabularx}
\end{table}

\reportfigure{blind.pdf}{两份 AI v2 盲评的配对均差。括号为每位评阅者的可比题数。默认模式的两个切题判断方向一致，可读性有分歧。}{fig:blind}

两个评阅者对默认模式的切题配对均差分别为 $+1.67$ 和 $+1.50$ 分，可读性为 $+0.17$ 与 $-0.07$ 分，连贯性为 $+0.37$ 与 $+0.23$ 分，三维的配对数均为 30。意境的配对均差为 $+0.20$ 与 $+0.31$，两个分母分别为 25、26；这是两种输出均通过适用门槛的子集。两份切题评分的增长提供了语义层面的补充证据，贴近题词和自然表达的关系仍须通过实例及新的真人盲评判断。两个 AI 上下文可能共享判断偏好，分值用于描述这批输出。

\begin{poembox}
\small\noindent\textbf{同题对读}\quad 题词「{\versefont 孤舟、春雨、煙波}」，种子 11\par
\noindent 自由生成\quad {\versefont 春潮平地脈，曉氣入雲深。獨立蒼茫裏，無人識此心。}\par
\noindent 默认模式\quad {\versefont 煙波春渺渺，寂寂孤舟閑。春雨夜來雨，空江寒葦山。}
\end{poembox}

这组输出说明两个指标各有独立用途。自由模式用春潮、晓气和独立远眺写出了季候与心境，三个原词均未完整出现。默认模式为读者提供明确的烟波与孤舟，三个原词均可定位，「春雨夜來雨」读来较生硬。只统计原词无法解释这种分歧。系统需要更多让多个词自然共现的训练配对，以及能在新题上检验衔接和表达的评分。

\subsection{新实现与新题上的隔位重字实验}

原排序已对整句及重复双字扣分，仍可能留下同一句内隔着其他字再次出现的单字。新排序把每句相邻同字合并成一个连续段，再统计该字曾出现过多少次。因而「渺渺」保留为一个连续段，「春雨夜來雨」计一次隔位返回。设原分数为 $s(y)$，计数为 $r(y)$，新排序使用 $s(y)-0.18\,r(y)$。权重在新题生成前固定。两种排序按相同约束、候选布局和随机种子重建同一候选池，并只改变排序规则。CLI 与本地网页均可切换。

实际补测分为两部分。先在已有 90 道验证题（种子 11）检查排序，原排序与新排序的重复分别为 12/90 和 7/90。继而冻结 36 道新题，一、二、三词各 12 道，所有词均通过字表检查，词集合和原主测试互异。每题各运行三个种子。四候选原排序和新排序有 4 条输出变化，新题中的隔位重字由 13/108 降到 9/108。按题作为抽样单位的配对 bootstrap，率差为 $-3.70$ 个百分点，95\% 分位区间为 $[-7.41,-0.93]$ 个百分点。该结果对应被直接优化的重字指标。

\reportfigure{ranking.pdf}{补充实验。A 固定候选后比较两种排序的隔位重字率。B 为新题 GPU 热推理延迟，每个模式 108 次请求。每条线连接中位数和第 95 百分位数。}{fig:ranking}

新题上单候选、四候选原排序和四候选新排序均是 108/108 三项达标。RTX 4080 SUPER 的中位耗时分别为 0.059、0.255、0.262 秒，同句隔位重字分别为 15/108、13/108、9/108。单候选适合更重视等待时间的演示。两种四候选模式的时延接近。自由生成在两词、三词各 36 次仍无全词覆盖。补测支持一个小范围的工程结论，即在现有候选中可少选几首隔位重字明显的诗。个别改善后文本依然会有生硬词组，且本轮新的 AI 盲评用的是原默认排序。新排序的语义与意境仍待同题盲评。

\section{部署、演示与复现}
\label{sec:deployment}

\subsection{交互与运行方式}

展示页让读者在同一个位置比较两首原文。输入限制为一至三个互异的短词，每词一至二字。界面同时展示偶句韵脚、韵部、字面命中与耗时，自动指标沿用冻结测试的口径。静态页面能够直接查看四组封存样例，现场演示可启动回环接口接收自选词。Python 进程启动时载入一次模型和字表，后续用同一模型顺序处理两种解码。参数合法性和字表检查在服务端再次执行，字表外字符、重复词等输入会返回提示。

\reportfigure{demo.png}{Orchid 展示页的实际浏览器画面。同题「霜、孤舟」，自由生成的「閑／山」与默认生成的「尋／林」都能确认同韵。默认生成同时含两个原词。}{fig:demo}

默认四候选模式在本机 CPU、两线程下的常驻中位耗时为 2.532 秒，在 RTX 4080 SUPER 为 0.198 秒，PyTorch 的峰值分配约为 370 MiB。相同权重的自由模式分别为 0.622 和 0.050 秒。这组速度实验每种模式、每设备预热两次再测 30 次，计时范围包含生成和排序，载入权重另测约 0.8 至 1.2 秒。展示接口在多次生成间保留模型。权重文件、字表和配置配套校验，CPU 可以只安装 CPU 构建依赖。

\reportfigure{latency.pdf}{CPU 与 GPU 上的生成时延。模型常驻，计时包含候选生成及排序。}{fig:latency}

\begin{tcolorbox}[consolestyle]
\begin{Verbatim}[fontsize=\small,breaklines=true]
python -m venv .venv
source .venv/bin/activate
python -m pip install -r requirements-cpu.txt
python -E scripts/serve_demo.py --checkpoint model --device cpu
\end{Verbatim}
\end{tcolorbox}

上面的命令在解压的模型包根目录运行，随后浏览 \url{http://127.0.0.1:8765/}。仓库开发目录使用 \code{dist/orchid-92m-prosody} 作为检查点路径。离线样例入口为 \code{demo/index.html}。

在实时页面可切换新的候选排序。命令行使用 \code{--ranking repeat-aware}。测试涵盖有效输入的双诗生成、字表外输入和错误词数。最终推理权重指纹以 \code{1613b489754c} 开头，完整值在模型包的 \code{bundle.json}。

\subsection{复现路径与验收}

公开源码仓库为 \url{https://github.com/DannyWANGD/Orchid}。可\href{https://codeload.github.com/DannyWANGD/Orchid/zip/040360649f565f08109eadd4bb861219cb63256b}{下载实验版本 \texttt{0403606} 的源码 ZIP}。解压后先阅读 \path{README.md}，模型实现位于 \path{orchid/model.py}，预训练和微调入口分别是 \path{scripts/pretrain.py} 和 \path{scripts/finetune.py}。配套推理包文件名为 \path{Orchid-model-kit-demo.zip}。

复现分为生成已训练模型、重跑冻结评估及从头训练。第一条路径用模型包内权重和离线韵表即可试用，不需要原始唐宋诗。第二条路径先按仓库说明从固定清单构建数据和监督集，以供应训练重合检查；然后将推理包传给 \code{evaluate_prosody.py}。脚本确认推理包沿袭已选检查点，并检查验证与训练文件指纹。用同一 CUDA 条件逐条比对了 600 条冻结输出，又在独立目录重新下载 313 个分片，核对了两套语料、划分、字表和监督统计。第三条路径从语料开始运行 \code{pretrain.py}、\code{finetune.py} 和韵律续训，以本次训练的验证集合选择自己的权重。种子固定能方便同设备比对，CPU 与 CUDA 的采样文本可能有差异。

\begin{tcolorbox}[consolestyle]
\begin{Verbatim}[fontsize=\small,breaklines=true]
python -E scripts/generate.py --checkpoint model --device cpu \
  --keywords 霜 孤舟 --seed 11 --ranking repeat-aware --json
python -E scripts/evaluate_ranking.py --split challenge \
  --checkpoint model --device cuda --output runs/ranking-replay
python -E -m unittest discover -s tests -v
\end{Verbatim}
\end{tcolorbox}

工程整理将训练与绘图的 Python 依赖压到 42 个固定项，单独提供 13 项的 CPU 推理依赖。发布源码保留模型实现、冻结配置、数据获取命令、图表生成代码和汇总证据。组织者盲评表与回件保存在本机。旧零回件状态与早期演示快照已清理，正式训练日志与指纹继续用于比较，使同一结果能够从数值追溯到模型和原文。

\begin{table}[!htbp]
\centering
\caption{复现所需的输入和入口。}
\label{tab:deliverables}
\begin{tabularx}{\linewidth}{@{}p{.20\linewidth}Y Y@{}}
\toprule
使用目的 & 已整理材料 & 核对方法\\
\midrule
查看示例 & \texttt{demo/index.html} 与四组离线封存原诗 & 同题原文、韵脚、种子与耗时可见\\
运行推理 & 92M 模型包、字表、韵表与 CPU/CUDA 依赖 & 核对 bundle 哈希，命令行与本地网页试用\\
重跑实验 & 源码、配置、冻结诗题、训练日志、评分说明和汇总 & 对照机器可读记录及新排序逐条生成\\
从头训练 & 固定公开唐宋诗版本与数据准备脚本，见参考资料 & 校验语料、分割、字表及监督数据指纹\\
重建报告 & LaTeX 源码、矢量图、浏览器截图和图表脚本 & XeLaTeX 编译并检查目录、交叉引用与字体\\
\bottomrule
\end{tabularx}
\end{table}

\section{总结、局限与项目体会}
\label{sec:conclusion}

本项目完成了从零训练的字符级 GPT、面向古诗的监督微调、韵律对照和可交互演示。模型扩容使固定语料上的验证损失下降，扩展唐宋语料进一步改善了同一留出集的预测。训练端的同韵样本续训增加了自由生成的偶句押韵。部署端的规则掩码保证了被测范围内的句式、押韵和字面用词。通过独立报告这两类结果，可以辨认哪些改善来自权重、哪些来自可控的生成过程。

最显著的限制来自多词条件生成。即使采用 92M 权重和同韵续训，自由模式在两套测试的两词、三词输入上都没有覆盖全词。按提示落点约束会带来重复、语序牵强和意象断裂，重字排序只能处理其中一种。改进训练时，可从同一首诗中派生语义共现的多意象样本，配平一词、两词和三词题，保留同一预算的对照。下一轮评估应冻结另一批题，由真人独立审阅语义切题和表达，并保留不同版本评分的适用数。格律方面可单列严格式，无需把每首四句诗都标成合乎近体常格的绝句。

课程实践中最有价值的体会是，评价目标需要和模型能力对应。预训练损失衡量对留出诗文的预测，监督损失衡量给定条件后对参考正文的拟合，写作任务还需要看真实抽样。只观察损失下降会错过两词与三词的明显退化，只看规则通过率也会错过一首诗读起来是否自然。项目中的人评反馈推动了 v2 量表，把字面命中和语义切题拆开，给连贯性和意境规定判定前提。这一调整和模型本身的改进同样重要。

从工程角度，检查点记录了数据、权重和游标，使长时间训练可以中断后继续，也让评估结果能追溯到一个确切权重。单卡试验表明显存占用与有效吞吐并不同时提高。真正的复现检查来自在新环境重新安装依赖、构建数据、解压运行以及重放逐条输出。项目作者确定古诗方向、完成早期读者评阅并提出格律与量表修订。代码实现、实验试跑和报告整理使用了 AI 协助，文中 AI 评阅有独立标注。教学实现查阅了 Raschka 的课程仓库和 nanoGPT~\cite{rasbt,nanogpt}，模型结构与训练脚本均能在附件中核对。


\clearpage
\phantomsection
\addcontentsline{toc}{section}{参考资料}
\begingroup
\footnotesize\linespread{1.04}\selectfont\raggedright
\setlength{\parskip}{.17em}
\bibliographystyle{unsrturl}
\bibliography{references}
\endgroup
\end{document}
